%!TEX program = xelatex
\documentclass[dvipsnames, svgnames,a4paper,11pt]{article}
% ----------------------------------------------------
%   中山大学物理与天文学院本科实验报告模板
%   作者：Huanyu Shi，2019级
%   知乎：https://www.zhihu.com/people/za-ran-zhu-fu-liu-xing
%   Github：https://github.com/huanyushi/SYSU-SPA-Labreport-Template
%   Last update : 2023.4.10
% ----------------------------------------------------

\input{Settings} % 导入模板的相关设置
\usepackage{lipsum}
\usepackage{enumitem}
\usepackage{tabularray}  %绘制表格时可以更加方便添加框线
\usepackage{longtable}
\setlist[enumerate]{label=\textup{(\arabic*)}}



%---------------------------------------------------------------------
%	正文
%---------------------------------------------------------------------

\begin{document}


\begin{table}
	\renewcommand\arraystretch{1.7}
	\begin{tabularx}{\textwidth}{
		|X|X|X|X
		|X|X|X|X|}
	\hline
	\multicolumn{2}{|c|}{预习报告}&\multicolumn{2}{|c|}{实验记录}&\multicolumn{2}{|c|}{分析讨论}&\multicolumn{2}{|c|}{总成绩}\\
	\hline
	\LARGE25 & & \LARGE30 & & \LARGE25 & & \LARGE80 & \\
	\hline
	\end{tabularx}
\end{table}


\begin{table}
	\renewcommand\arraystretch{1.7}
	\begin{tabularx}{\textwidth}{|X|X|X|X|}
	\hline
	专业：& 物理学 &年级：& 2022级\\
	\hline
	姓名：& 戴鹏辉 \& 杨舒云  & 学号： & 2344016 \& 22344020 \\
	\hline
	日期：& 2024/11/5 & 教师签名：& \\
	\hline
	\end{tabularx}
\end{table}

\begin{center}
	\LARGE D6 \quad $\gamma$能谱测量实验
\end{center}

\textbf{【实验报告注意事项】}
\begin{enumerate}
	\item 实验报告由三部分组成：
	\begin{enumerate}
		\item 预习报告：（提前一周）认真研读\underline{\textbf{实验讲义}}，弄清实验原理；实验所需的仪器设备、用具及其使用（强烈建议到实验室预习），完成课前预习思考题；了解实验需要测量的物理量，并根据要求提前准备实验记录表格（第一循环实验已由教师提供模板，可以打印）。预习成绩低于10分（共20分）者不能做实验。
	    \item 实验记录：认真、客观记录实验条件、实验过程中的现象以及数据。实验记录请用珠笔或者钢笔书写并签名（\textcolor{red}{\textbf{用铅笔记录的被认为无效}}）。\textcolor{red}{\textbf{保持原始记录，包括写错删除部分，如因误记需要修改记录，必须按规范修改。}}（不得输入电脑打印，但可扫描手记后打印扫描件）；离开前请实验教师检查记录并签名。
	    \item 分析讨论：处理实验原始数据（学习仪器使用类型的实验除外），对数据的可靠性和合理性进行分析；按规范呈现数据和结果（图、表），包括数据、图表按顺序编号及其引用；分析物理现象（含回答实验思考题，写出问题思考过程，必要时按规范引用数据）；最后得出结论。
	\end{enumerate}
	\textbf{实验报告就是将预习报告、实验记录、和数据处理与分析合起来，加上本页封面。}
	% \item 每次完成实验后的一周内交\textbf{实验报告}（特殊情况不能超过两周）。
	\item 实验安全注意事项
		\begin{enumerate}[label=\roman*.]
			\item 使用真实放射源进行实验时，使用镊子放置和移除放射源，实验完毕后将放射源放回源库，并洗手。
		\end{enumerate}
\end{enumerate}


\clearpage
\tableofcontents
\clearpage

\setcounter{section}{0}
\section{D6 \quad $\gamma$能谱测量实验 \quad\heiti 预习报告}
	
\subsection{实验目的}
\begin{enumerate}
	\item 测定 $\gamma$ 谱仪的能量分辨率以及能量线性；
	\item 测定已知源的 $\gamma$ 能谱，并作能谱分析；
	\item 随机数据处理；
	
\end{enumerate}

\subsection{仪器用具}


\begin{table}[htbp]
	\centering
	\renewcommand\arraystretch{1.6}
	% \setlength{\tabcolsep}{10mm}
	\caption{$\gamma$ 能谱测量实验用具}
	\begin{tabular}{p{0.05\textwidth}|p{0.20\textwidth}|p{0.05\textwidth}|p{0.5\textwidth}}
		\hline
		编号& 仪器用具名称 & 数量 &  主要参数（型号，测量范围，测量精度等） \\
		\hline
		1	&	核信号发生器 	& 1 & NMS-6014-SIG，输出信号电压范围 0-5V	\\

		2	&	通用核信息采集 	& 1 & NMS-6014-GD，采样频率 250MHz，带宽100MHz；输入范围 0-5V，8192 道，最大计数率10M/s，死时间 100ns；	 \\
		
		3	&	计算机 & 1 & 安装有 Labview2014 及虚拟核数据采集软件	 	\\
		
		4	&	示波器	& 1 &  RIGOL DS2202A，带宽 200MHZ，双通道，采样频率 1GSa/s	\\
		
		\hline
	\end{tabular}
	\label{tbl:gamma 能谱测量实验用具}
\end{table}




% \clearpage

\subsection{原理概述}

	伽马射线是一种波长极短（小于 0.2 \AA）的高能电磁波，具有显著的粒子特性、极高的能量和极强的穿透能力。伽马射线的来源包括粒子碰撞、正反粒子湮灭、能级跃迁和韧致辐射等过程。

	
	\begin{enumerate}
		\item \textbf{伽马射线与物质的相互作用：}

			\begin{itemize}
				\item \textbf{光电效应：}
				
					当能量$E_{\gamma}$ 的入射 $\gamma$ 光子与物质中原子的束缚电子相互作用时，光子可以把全部能量转移给某个束缚电子，使电子脱离原子束缚而发射出去，光子本身消失，发射出去的电子称为光电子，这种过程称为光电效应。发射出光电子的动能
					\[
						E_e = E_{\gamma} - B_i
					\]
					$B_i$ 为束缚电子所在壳层的结合能。
		
		
				\item \textbf{康普顿效应：}
				
					$\gamma$ 光子与自由静止的电子发生碰撞，而将一部分能量转移给电子，使电子成为反冲电子，$\gamma$ 光子被散射改变了原来的能量和方向，反冲电子的动能为：
					\[
						E_e = \frac{E_{\gamma}^2 (1 - \cos\theta)}{m_0 c^2 + E_{\gamma}(1 - \cos\theta)} = \frac{E_{\gamma}}{1 + \frac{m_0 c^2}{E_{\gamma} (1 - \cos\theta)}}
					\]
					$m_0 c^2$ 为电子静止质量，角度 $\theta$ 是 $\gamma$ 光子的散射角。当$\theta = 180^\circ$ 时，反冲电子的动能有最大值：
					\[
						E_{e,max} = \frac{E_\gamma}{1 + \frac{m_0 c^2}{2 E_\gamma}}
					\]
					这说明康普顿效应产生的反冲电子的能量有一上限最大值，称为康普顿边界。
					
					反冲电子以角度 $\varphi$ 出射，则$\varphi$ 和 $\theta$ 存在关系：
					\[
						\cot\varphi = (1 + \frac{E_{\gamma}}{m_0 c^2})\tan\frac{\theta}{2}
					\]
		
		
		
				\item \textbf{电子对效应：}
		
					当 $\gamma$ 光子能量大于 $2 m_0 c^2$ 时，$\gamma$ 光子从原子核旁边经过并受到核的库仑场作用，可能转化为一个正电子和一个负电子，称为电子对效应。此时光子能量可表示为两个电子的动能与静止能量之和：
					\[
						E_\gamma = E_e^+ + E_e^- + 2 m_0 c^2
					\]
		
		
		
		
		
		
			\end{itemize}

		\item \textbf{谱仪的能量刻度和分辨率：}
			
			应用 $\gamma$ 谱仪测定未知射线能谱时，必须先用已知能量的核素能谱来标定 $\gamma$ 谱仪。由于能量与各峰位道址是线性的：$E = k \cdot d + b$，因此能量刻度就是设法得到 $k$ 和 $b$。

			$\gamma$ 能谱仪的一个重要指标是能量分辨率。定义谱仪能量分辨率：
			\[
				\eta = \frac{FWHM}{E_\gamma} \times 100 \%
			\]
			其中 FWHM(Full Width Half Maximum)表示选定能谱峰的半高全宽，$E_\gamma$为与谱峰对应的$\gamma$光子能量，$\eta$表示闪烁谱仪在测量能量时能够分辨两条靠近的谱线的本领。
			
	\end{enumerate}




\subsection{实验前思考题}

% 思考题1
\begin{question}
	简述 $\gamma$ 光子与物质的相互作用类型，及其产生的次级粒子与原入射光子的能量关系。
\end{question}

	% \begin{itemize}
	% 	\item \textbf{光电效应：}
		
	% 		当能量$E_{\gamma}$ 的入射 $\gamma$ 光子与物质中原子的束缚电子相互作用时，光子可以把全部能量转移给某个束缚电子，使电子脱离原子束缚而发射出去，光子本身消失，发射出去的电子称为光电子，这种过程称为光电效应。发射出光电子的动能
	% 		\[
	% 			E_e = E_{\gamma} - B_i
	% 		\]
	% 		$B_i$ 为束缚电子所在壳层的结合能。


	% 	\item \textbf{康普顿效应：}
		
	% 		$\gamma$ 光子与自由静止的电子发生碰撞，而将一部分能量转移给电子，使电子成为反冲电子，$\gamma$ 光子被散射改变了原来的能量和方向，反冲电子的动能为：
	% 		\[
	% 			E_e = \frac{E_{\gamma}^2 (1 - \cos\theta)}{m_0 c^2 + E_{\gamma}(1 - \cos\theta)} = \frac{E_{\gamma}}{1 + \frac{m_0 c^2}{E_{\gamma} (1 - \cos\theta)}}
	% 		\]
	% 		$m_0 c^2$ 为电子静止质量，角度 $\theta$ 是 $\gamma$ 光子的散射角。当$\theta = 180^\circ$ 时，反冲电子的动能有最大值：
	% 		\[
	% 			E_{e,max} = \frac{E_\gamma}{1 + \frac{m_0 c^2}{2 E_\gamma}}
	% 		\]
	% 		这说明康普顿效应产生的反冲电子的能量有一上限最大值，称为康普顿边界。
			
	% 		反冲电子以角度 $\varphi$ 出射，则$\varphi$ 和 $\theta$ 存在关系：
	% 		\[
	% 			\cot\varphi = (1 + \frac{E_{\gamma}}{m_0 c^2})\tan\frac{\theta}{2}
	% 		\]



	% 	\item \textbf{电子对效应：}

	% 		当 $\gamma$ 光子能量大于 $2 m_0 c^2$ 时，$\gamma$ 光子从原子核旁边经过并受到核的库仑场作用，可能转化为一个正电子和一个负电子，称为电子对效应。此时光子能量可表示为两个电子的动能与静止能量之和：
	% 		\[
	% 			E_\gamma = E_e^+ + E_e^- + 2 m_0 c^2
	% 		\]






	% \end{itemize}

	详见上述实验原理部分。
	





% 思考题2
\begin{question}
	对探测器中的光电倍增管，射线的光电效应会对其测量产生影响吗？影响能有多大？
\end{question}

	

	\begin{enumerate}
		\item 光电子的产生：当光照射到光电倍增管的光阴极时，根据光电效应，光阴极会激发出光电子。这些光电子的数量直接受到入射光强度的影响，进而影响光电倍增管的信号强度。
		\item 信号的放大：产生的光电子通过电子光学系统进入倍增系统，在倍增系统中，每个光电子通过二次发射效应被放大，最终在阳极形成较大的电流或电压输出。因此，光电效应的效率直接影响了光电倍增管的灵敏度和信号放大能力。
		\item 时间分辨率：光电倍增管的时间分辨率受到光电子信号在各个部分传输时间的影响，包括从光阴极到微通道板的传输时间、在微通道板中的传输时间、从微通道板到阳极的传输时间，以及在阳极到电极端口的传输时间。光电效应中光电子的产生和传输速度对时间分辨率有直接影响。
		\item 测量误差：不同位置的光电子信号在阳极上会产生不同的测量误差，这种误差会影响到信号到达阳极的时间测量，从而影响系统的时间分辨率。因此，光电效应中光电子的产生位置和传输路径对测量结果的准确性有重要影响。
		\item 噪声和信号质量：光电效应还可能引入噪声，如散粒噪声、热噪声等，这些噪声会影响光电倍增管测量的信号质量。
	\end{enumerate}
	

	综上所述，光电效应对光电倍增管的测量有显著影响，包括信号的产生、放大、时间分辨率以及测量误差等多个方面。为了提高测量的准确性和时间分辨率，需要对光电效应中产生的光电子进行精确控制和优化，以减少误差和噪声的影响。


	




% 思考题3
\begin{question}
	有一单能 $\gamma$ 源，能量为 2MeV，根据 $\gamma$ 与物质的相互作用及 NaI（TI）闪烁能谱仪输出信号幅度的关系，预测能谱形状。
\end{question}

	\begin{itemize}
		\item 光电峰（Photopeak）：在$\gamma$射线与NaI（TI）晶体相互作用时，光电效应是主要的过程之一。在光电效应中，$\gamma$光子被完全吸收，将其全部能量转移给一个电子，使该电子从原子中逸出，称为光电子。这个过程中，$\gamma$射线的能量完全转化为光电子的能量，因此，在能谱上会形成一个对应于$\gamma$射线能量的峰，即光电峰。对于2MeV的$\gamma$源，我们可以预期在2MeV处会有一个明显的峰。
		\item 康普顿峰（Compton Edge）和康普顿连续谱（Compton Continuum）：康普顿效应是$\gamma$射线与物质相互作用的另一个重要过程。在康普顿散射中，$\gamma$光子与自由电子或外层电子相互作用，将部分能量转移给电子，散射出一个能量较低的$\gamma$光子和一个康普顿电子。康普顿散射后，散射光子的能量与散射角有关，因此，康普顿效应会产生一个连续的能量谱，从零能量开始，直到康普顿边（Compton Edge），即散射光子能量等于入射$\gamma$射线能量减去康普顿电子最大能量时的能量。对于2MeV的$\gamma$源，康普顿边的能量可以通过康普顿散射公式计算得出，大约在1.022MeV（2MeV - 2$m_ec^2$，其中$m_ec^2$是电子静止质量对应的能量）。
		\item 全能峰（Total Absorption Peak）：全能峰是指γ射线在晶体中被完全吸收，包括光电效应和康普顿效应后，康普顿散射产生的散射光子再次被晶体吸收，通过光电效应转换成光电子能量，这些光电子对输出脉冲的贡献叠加在一起形成一个脉冲。这个脉冲幅度所对应的能量，是这两个电子的能量之和，即$E_e + h \nu' = h \nu$，即等于入射γ射线的能量。因此，全能峰将叠加在光电峰之上，使之增高。
	\end{itemize}
	

	综上所述，对于一个单能γ源，能量为2MeV，我们可以预期在能谱上看到一个明显的光电峰在2MeV处，一个康普顿连续谱从低能量开始，直到康普顿边（约1.022MeV），以及全能峰叠加在光电峰之上。这样的能谱形状可以帮助我们理解γ射线与物质相互作用的特性，并用于γ射线能量的测量和分析
。







% 思考题4
\begin{question}
	闪烁谱仪是常用的 $\gamma$ 能谱仪，简述闪烁谱仪的的结构图和各部分功能。
\end{question}

	\begin{itemize}
		\item 闪烁体：闪烁体是闪烁谱仪的核心部分，通常由NaI(Tl)晶体构成。当γ射线入射至闪烁体时，会通过光电效应、康普顿效应和电子对效应与闪烁体发生相互作用，产生次级电子，这些电子使闪烁体分子电离和激发，退激时发出光子。
		\item 光电倍增管（PMT）：光电倍增管位于闪烁体的附近，用于探测闪烁体发出的光子。光子在光电倍增管的光阴极上打出光电子，这些光电子在光电倍增管中倍增，形成大量的电子，最终被阳极接收形成电压脉冲。
		\item 高压电源：高压电源为光电倍增管提供所需的高压，以维持光电子的倍增过程。
		\item 线性放大器：线性放大器用于放大从光电倍增管传来的电压脉冲，并进行成形处理，以便后续的脉冲幅度分析。
		\item 多道脉冲幅度分析器：多道脉冲幅度分析器（或单道脉冲幅度分析器）用于记录不同幅度的脉冲，通过分析脉冲幅度可以确定入射γ射线的能量。
		\item 定标器：定标器用于校准脉冲幅度分析器，确保测量结果的准确性。
		\item 射极跟随器：射极跟随器用于阻抗匹配，将光电倍增管产生的电压脉冲传输到线性脉冲放大器。
		\item 反射层：在闪烁体与光电倍增管之间通常会有一层反射层，用于反射闪烁体发出的光子，以提高光收集效率。
		\item 计算机系统：用于控制谱仪的操作，以及对采集到的数据进行处理和分析。
	\end{itemize}
	

	闪烁谱仪的工作原理是：γ射线与闪烁体相互作用产生光子，这些光子被光电倍增管探测并转换成电信号，然后通过电子学设备进行放大、分析和记录，最终得到γ射线的能谱。通过测量不同能量γ射线产生的脉冲幅度分布，可以对γ射线的能量进行分析和测量
。






\clearpage
\begin{table}
	\renewcommand\arraystretch{1.7}
	\centering
	\begin{tabularx}{\textwidth}{|X|X|X|X|}
	\hline
	专业：& 物理学 &年级：& 2022级 \\
	\hline
	姓名：& 戴鹏辉 \& 杨舒云 & 学号：& 22344016 \& 22344020 \\
	\hline
	室温：& 26℃ & 实验地点： & A404 \\
	\hline
	学生签名：& 见附件 & 评分： &\\
	\hline
	实验时间：& 2024/10/31& 教师签名：&\\
	\hline
	\end{tabularx}
\end{table}

\section{D6 \quad $\gamma$ 能谱测量实验 \quad\heiti 实验记录}

\subsection{实验内容和步骤}

	\begin{itemize}
		\item \textbf{理论计算$^{137}\mathrm{Cs}$康普顿散射峰和背散射峰：}

			\begin{enumerate}
				\item 已知$^{137}\mathrm{Cs}$主峰能量为：$E_{\gamma} = 0.662 \ \mathrm{(MeV)}$。
				\item 理论计算得到康普顿散射峰能量：$E_{max} = \frac{E_{\gamma}}{1 + \frac{m_0 c^2}{2 E_{\gamma}}} = 0.47765 \ \mathrm{(MeV)}$
				\item 背散射峰：$E_b = E_{\gamma} - E_{max} = 0.662 - 0.47765 \ \mathrm{(MeV)} = 0.18345 \ \mathrm{MeV}$
			\end{enumerate}


		\item \textbf{测量$^{137}\mathrm{Cs}$和$^{60}\mathrm{Co}$的能谱图：}
			
			测量二者的能谱图，得到结果如\cref{fig:Data1-主峰}、\cref{fig:Data3-主峰1}：

			\begin{figure}[htbp]
				\centering
				\includegraphics[width=0.7\linewidth]{Data1-主峰.png}
				\caption{$^{137}\mathrm{Cs}$的能谱图}
				\label{fig:Data1-主峰}
			\end{figure}

			\begin{figure}[htbp]
				\centering
				\includegraphics[width=0.7\linewidth]{Data3-主峰1.png}
				\caption{$^{60}\mathrm{Co}$的能谱图}
				\label{fig:Data3-主峰1}
			\end{figure}

			手动寻峰得到的测量结果为：
			\begin{itemize}
				\item $^{137}\mathrm{Cs}$放射源：主峰峰位：431
				\item $^{137}\mathrm{Cs}$放射源：康普顿散射峰峰位：289
				\item $^{137}\mathrm{Cs}$放射源：背散射峰峰位：141
				\item $^{60}\mathrm{C0}$放射源：主峰1峰位：865
				\item $^{60}\mathrm{C0}$放射源：主峰2峰位：771
			\end{itemize}



		\item \textbf{估算能量分辨率：}
		
			由\cref{fig:Data1-主峰}可得，内置软件的计算结果为：

			\begin{itemize}
				\item 峰位：420.33
				\item 半高宽：45.34
				\item 半高宽/峰位：10.79\%
			\end{itemize}

			但是由于肉眼可见的寻峰有误，这个结果肯定也是有误的。我们将在后面的数据分析中解决这个问题。
			
			% 对所得到的能谱进行寻峰、并获取全高半宽，得到结果如\cref{fig:D6-3-2}：
			
			% \begin{figure}[htbp]
			% 	\centering
			% 	\includegraphics[width=0.7\linewidth]{D6-3-2.png}
			% 	\caption{$^{137}\mathrm{Cs}$主峰与全高半宽}
			% 	\label{fig:D6-3-2}
			% \end{figure}

			% \begin{enumerate}
			% 	\item $^{137}\mathrm{Cs}$放射源：主峰峰位：432.0
			% 	\item $^{137}\mathrm{Cs}$放射源：主峰峰位半高宽：47.28
			% \end{enumerate}

			% 于是可计算多道分析仪能量分辨率：
			% \begin{align*}
			% 	\eta = \frac{47.28 \times }{}
			% \end{align*}

	\end{itemize}
		




		


	\subsection{实验过程中遇到的问题记录}

		\begin{enumerate}
			\item 实验中测量$^{60}\mathrm{Co}$的能谱图时，发现两个主峰非常不明显。可能的原因是“模拟软件的板卡数据有误”，但是主峰的位置是正确的，不影响数据分析。
			\item 由于软件内置的寻峰算法效果较差，肉眼可见的寻峰有误，所以只能进行手动寻峰，人工读数；在半高宽的测量值，人工读数也不可靠，只能通过后期的数据分析解决。
				
		\end{enumerate}



\subsection{实验数据记录}

	见\cref{fig:Data}

	\begin{figure}[htbp]
		\centering
		\includegraphics[width=0.7\linewidth]{Data-gamma.jpg}
		\caption{原始数据记录}
		\label{fig:Data}
	\end{figure}

	% \begin{figure}[htbp]
	% 	\centering
	% 	% 第一行的三张图片
	% 	\subfloat[原始数据1]
	% 	{\includegraphics[width=0.3\textwidth]{D3-OriginalData-1.jpg}\label{fig:data1}}
	% 	\quad
	% 	\subfloat[原始数据2]
	% 	{\includegraphics[width=0.3\textwidth]{D3-OriginalData-2.jpg}\label{fig:data2}}
	% 	\quad
	% 	\subfloat[原始数据3]
	% 	{\includegraphics[width=0.3\textwidth]{D3-OriginalData-3.jpg}\label{fig:data3}}
	% 	\newline
	% 	% 第二行的两张图片
	% 	\subfloat[原始数据4]
	% 	{\includegraphics[width=0.3\textwidth]{D3-OriginalData-4.jpg}\label{fig:data4}}
	% 	\quad
	% 	\subfloat[原始数据5]
	% 	{\includegraphics[width=0.3\textwidth]{D3-OriginalData-5.jpg}\label{fig:data5}}
	% 	\caption{原始数据}
	% 	\label{fig:data}
	% \end{figure}
	



%\subsection{原始数据记录}


% \clearpage

	

\clearpage
\begin{table}
	\renewcommand\arraystretch{1.7}
	\begin{tabularx}{\textwidth}{|X|X|X|X|}
	\hline
	专业：& 物理学 &年级：& 2022级\\
	\hline
	姓名： & 戴鹏辉 \& 杨舒云 & 学号：& 22344016 \& 22344020\\
	\hline
    日期：& 2024/11/15 & 评分： &\\
	\hline
	\end{tabularx}
\end{table}

\section{D6 \quad $\gamma$ 能谱测量实验 \quad\heiti 分析与讨论}

\subsection{实验数据分析}


	\subsubsection{对探测器进行能量刻度标定：}

		由于软件内置的寻峰算法效果较差，我们重新使用了自己的寻峰算法，对数据进行了如下处理：
		\begin{itemize}
			\item 使用 Savitzky-Golay 滤波器对数据进行平滑，以减少噪声对主峰检测的影响。
			\item 使用 scipy. signal. find\_peaks 函数找到信号中的所有峰值，并确定主峰（最高的峰值）。
			\item 计算主峰的半高宽。
			\item 结果如\cref{fig:D6-3-2}、\cref{fig:D6-3-3}：
		\end{itemize}


		\begin{figure}[htbp]
			\centering
			\includegraphics[width=0.7\linewidth]{D6-3-2.png}
			\caption{$^{137}\mathrm{Cs}$主峰与全高半宽}
			\label{fig:D6-3-2}
		\end{figure}

		\begin{figure}[htbp]
			\centering
			\includegraphics[width=0.7\linewidth]{D6-3-3.png}
			\caption{$^{60}\mathrm{Co}$滤波后的寻峰结果}
			\label{fig:D6-3-3}
		\end{figure}





		\begin{enumerate}
			\item 使用三个已知能量的峰位进行能量刻度标定：
				\begin{itemize}
					\item $^{137}\mathrm{Cs}$放射源：主峰峰位：432.0，对应能量：0.662MeV
					\item $^{60}\mathrm{Co}$放射源：主峰1峰位：858，对应能量：1.332MeV
					\item $^{60}\mathrm{Co}$放射源：主峰2峰位：755，对应能量：1.173MeV
				\end{itemize}

			\item 由测量结果，对探测器进行能量标定，得到结果如\cref{fig:D6-3-1}：
				
				\begin{figure}[htbp]
					\centering
					\includegraphics[width=0.7\linewidth]{D6-3-1.png}
					\caption{对探测器进行能量标定}
					\label{fig:D6-3-1}
				\end{figure}


			\item 得到拟合结果为：
			\[
				E = 0.0016 \cdot d - 0.0180 \ (\mathrm{Me V})
			\]

		\end{enumerate}


	\subsubsection{计算$^{137}\mathrm{Cs}$的康普顿散射峰和背散射峰能量与测量值的相对误差}

		\begin{enumerate}
			\item 理论计算得到：
				\begin{itemize}
					\item 康普顿散射峰能量：$E_{max} = 0.47765 \ \mathrm{(MeV)}$
					\item 康普顿背散射峰能量：$E_b = E_{\gamma} - E_{max} = 0.18345 \ \mathrm{(MeV)}$
				\end{itemize}

			\item 实验测量得到：
				\begin{enumerate}
					\item 康普顿散射峰峰位：289 $\Rightarrow E'_{max} = 0.0016 \times 289 - 0.0180 \ \mathrm{(MeV)} = 0.4444 \ \mathrm{(MeV)}$
					\item 康普顿背散射峰峰位：141 $\Rightarrow E'_b = 0.0016 \times 141 - 0.0180 \ \mathrm{(MeV)} = 0.2076 \ \mathrm{(MeV)}$
				\end{enumerate}

			\item 实验值与理论值的相对误差为：
				\begin{enumerate}
					\item $\eta_{E_{max}} = \frac{0.4444 - 0.47765}{0.47765} \times 100 \% =  -6.961 \% $
					\item $\eta_{E_b} = \frac{0.2076 - 0.18345}{0.18345} \times 100 \% = 13.164 \% $
				\end{enumerate}

			
			\item 可以看到，康普顿散射峰与背散射峰的实验值与理论值之间的相对误差均在15 \% 之内。误差的来源可能是“康普顿散射峰与背散射峰的实验时的道址数”不准确，这两个峰并不在能谱的极大值位置，所以难以使用寻峰函数直接寻找，只能手动读数，导致误差较大。 


		\end{enumerate}





	\subsubsection{估算多道分析仪能量分辨率}

		由\cref{fig:D6-3-2}，有：

		\begin{enumerate}
			\item $^{137}\mathrm{Cs}$放射源：主峰峰位：432.0
			\item $^{137}\mathrm{Cs}$放射源：主峰峰位半高宽：47.28
		\end{enumerate}

		于是可计算多道分析仪能量分辨率：
		\begin{align*}
			\eta = \frac{47.28 \times 0.0016 - 0.0180 \ \mathrm{(MeV)}}{432.0 \times 0.0016 - 0.0180 \ \mathrm{(MeV)}} \times 100 \% = 8.563 \%
		\end{align*}


	% \subsubsection{实验四 \quad XXXXXXXXXXXXXXXXXXXXXXX}


		

	
% \subsection{实验后思考题}

% % 思考题1
% \begin{question}
% 	XXXXXXXXXXXXXXXXXXXXXXXXXXXXX
% \end{question}

	
% % 思考题2
% \begin{question}
% 	XXXXXXXXXXXXXXXXXXXXXX
% \end{question}


% % 思考题3
% \begin{question}
% 	XXXXXXXXXXXXXXXXXXXXXX
% \end{question}



\clearpage
\section{D6 \quad $\gamma$ 能谱测量实验 \quad\heiti 实验总结}

	\subsection{实验分工}

		本实验由戴鹏辉和杨舒云两人合作完成
		\begin{enumerate}
			\item 戴鹏辉完成实验操作部分和实验报告撰写部分。
			\item 杨舒云完成实验数据分析与处理部分。
		\end{enumerate}
	
		


	% \subsection{实验心得}

	\subsection{Attachment}
	% The arrangement of the experimental bench desktop is shown in %\cref{}.

	个人签名见 \cref{fig:name1} 、 \cref{fig:name2}.

	\begin{figure}[htbp]
		\centering
		\includegraphics[width=0.5\textwidth]{name.png}
		\caption{signature1}
		\label{fig:name1}
	\end{figure}
	
	\begin{figure}[htbp]
		\centering
		\includegraphics[width=0.35\textwidth]{name-TaLEs.jpg}
		\caption{signature2}
		\label{fig:name2}
	\end{figure}

	

	
	


\end{document}
